\section{n4. New all-support geometric cone and envelope invariance for finite and specified infinite Frobenius-root extensions.}
\textit{Supports: G331, n1, n3.}

NEW Frobenius-root invariance on the entire geometric envelope class. Let R subset S be an injective finite local extension of excellent reduced positive-dimensional rings of characteristic p\textgreater{}0. Assume a fixed q=p\textasciicircum{}e satisfies s\textasciicircum{}q in R for EVERY s in S. This hypothesis holds when R is a normal domain, S a finite domain, and Frac(S)/Frac(R) is finite purely inseparable: choose q killing that extension; s\textasciicircum{}q lies in Frac(R), is integral over R and hence lies in R. No perfection, F-finiteness, CM, Q-Gorenstein or resolution hypothesis is imposed. This is a proved envelope transport theorem; it supplies no unproved positive-characteristic regularity input.

Fix a primary ideal a in R. Restriction rho:V\_(aS)(S)-\textgreater{}V\_a(R) is an actual homeomorphism with inverse

 w(s)=v(s\textasciicircum{}q)/q, including infinity/q=infinity.

The formula has w|R=v. Frobenius additivity (s+t)\textasciicircum{}q=s\textasciicircum{}q+t\textasciicircum{}q and multiplicativity prove the semivaluation axioms. Because aS is primary and w(aS)=1, the source maximal ideal has positive finite value. For example n\textasciicircum{}N subset mS subset n gives 0\textless{}w(m)/N\textless{}=w(n)\textless{}=w(m)\textless{}infinity. Supports are determined by s\textasciicircum{}q in supp(v), including nonminimal supports. Any lift must have this value by multiplicativity, so it is unique. Both maps are continuous by coordinate evaluations and scalar division, or by compact bijection. This proves full fibre compatibility rather than assuming generic inseparability controls every support.

The independently geometric classes C and D are isometrically identical under rho. One may invoke G331's full finite-injection positive descent since every fibre is singleton, but there is also an explicit ideal proof. If I=(s\_1,...,s\_l) is unit or proper primary in S, put J=(s\_1\textasciicircum{}q,...,s\_l\textasciicircum{}q) in R. Then v(J)=q w(I) on EVERY support. If I contains (aS)\textasciicircum{}N, raising any expression x=sum\_i s\_i u\_i for x in a\textasciicircum{}N to the qth power shows x\textasciicircum{}q in J, since all u\_i\textasciicircum{}q are in R. Thus J has radical m; it is proper when I is proper and is unit when I is unit. Hence the rational profile -w(I)/r is exactly the pullback of -v(J)/(qr). Primary ideal extension gives the reverse inclusion. Rational density and uniform closure identify C; actual decreasing C-nets identify D in both directions. Thus for EVERY pointwise obstacle psi, whenever the raw supremum is defined with the usual bottom convention,

 rho\textasciicircum{}*p\_R(psi)=p\_S(rho\textasciicircum{}*psi).

For finite continuous obstacles, a continuous greatest geometric envelope exists on one side iff it exists on the other. In particular the all-obstacle CE property, any genuine continuous-obstacle CE failure, and continuity of the full canonical approximation transfer both ways. No norm average or positivity of a USC regularization is used. This is the exact operator compatibility required for these Frobenius-root maps, without any Monge-Ampere/product assertion.

There is quantitative canonical transport, even though same-degree scalar extension of canonical ideals is NOT asserted. Let B\_m=A\_m\textasciicircum{}S(rho\textasciicircum{}*f) and A\_m=A\_m\textasciicircum{}R(f). Surjectivity gives A\_m=B\_m intersect R. Let J\_m be generated by qth powers of generators of B\_m. Then kappa\_m=-v(J\_m)/(qm) pulls back to -w(B\_m)/m and J\_m subset A\_(qm). If a source estimate 0\textless{}=p\_S-h\_m\textasciicircum{}S\textless{}=E/m holds, n1's residue grading with M\_1=||v(A\_1)|| gives, for k\textgreater{}=2q,

 0\textless{}=p\_R-h\_k\textasciicircum{}R\textless{}=E/floor(k/q)+qM\_1/k\textless{}=(2qE+qM\_1)/k.

This is a conditional rate with its source hypothesis explicit. Uniform source convergence alone suffices for uniform target convergence, either by the exact envelope homeomorphism and n1 or by the same residue argument. Frobenius exponent q is not identified with generic degree or a ramification/residue sum.

An infinite-root version is also proved. Suppose R subset S is an injective extension of excellent reduced local rings of characteristic p, S is the union of local finite R-subalgebras, and every element of S has SOME p-power in R. Assume S is Noetherian and the marked maximal ideal and primary gauge are the extensions of those of R, as in the scalar example below. Define w(s)=v(s\textasciicircum{}(p\textasciicircum{}e))/p\textasciicircum{}e using any admissible e. Passing to a common larger exponent proves independence of e and proves all axioms for sums and products. Each coordinate formula uses a finite exponent, so the inverse is continuous; centredness follows from the extended maximal ideal. Every support is retained and rho is again a boundary homeomorphism. Every finitely generated primary source ideal and finitely many coefficients witnessing a primary-power containment have a common Frobenius exponent. Raising those expressions as above gives a genuine primary ideal J in R with exact rational profile. Thus C and actual D again coincide isometrically and envelopes commute for ALL obstacles. There need not be one global q, so no global degree-rate constant is claimed.

For a geometrically integral geometrically normal rational-point finite-type local k-domain in characteristic p, S=R tensor\_k k\textasciicircum{}perf is such an infinite-root example. It is the corresponding Noetherian excellent normal local ring, with mS its maximal ideal. Finite purely inseparable coefficient extensions form local finite stages. Each element is a finite sum of tensors; choose one exponent sending all its finitely many coefficients into k, and Frobenius sends the element into R. Thus CE is equivalent over k and its perfect closure, even for imperfect k, with every support included. This reduces the imperfect-field issue for this geometrically normal finite-type class; it does not prove CE in positive characteristic over the perfect closure.

Combining n3's separable-field descent with this root theorem gives a precise further reduction. For such geometrically normal rational-point germs in characteristic p, pass first to k\textasciicircum{}s and then to its perfect closure, which is an algebraic closure of k. Root invariance identifies the latter two envelope problems, and n3 gives rho\textasciicircum{}*p\_R(f)=p\_(R tensor k\textasciicircum{}s)(rho\textasciicircum{}*f) on compatible obstacles. Therefore geometric-algebraic-closure CE implies downstairs CE at these exact hypotheses, without enlarging BFJ to characteristic p. Mixed-characteristic rings have no Frobenius-root formula of this kind and are not covered. General nonnormal/conductor envelope commutation remains conditional outside these uniquely lifting maps, and all literal radial distinctions remain unchanged.

Artifacts: n4.txt, distinct full n4.tex, standalone TeX, entire companion and scoped verdict. Rerun: python assemble\_delivery.py; python finite\_checks.py; python check\_delivery.py. Finite Frobenius arithmetic is diagnostic only; the all-support cone and operator claims are proved above relative to qualified G331 and own n1/n3.

\subsection*{Expanded root transport and its quantitative limit}
For $q=p^e$ and $S^q\subset R$, all supports obey
\[
 w(s)=q^{-1}v(s^q),\qquad
 v((s_1^q,\ldots,s_l^q))=q\,w((s_1,\ldots,s_l)).
\]
The primaryness proof raises actual module expressions, not just
generic field identities: if $x\in\mathfrak a^N$ is expressed as
$\sum s_i u_i$, then $x^q=\sum s_i^qu_i^q\in J$ in $R$.
Thus the transported profiles have actual primary Cartier presentations.
Under an explicitly assumed source bound, and $k\ge2q$,
\[
 0\le p_R-h_k^R\le
 \frac{E}{\lfloor k/q\rfloor}+\frac{qM_1}{k}
 \le\frac{2qE+qM_1}{k}.
\]
An infinite coefficient perfection need not have bounded exponent.
Each finite ideal still has one common exponent, giving the exact
cone/operator transport, but the display supplies no uniform infinite
extension rate without additional bounds.

